% tikz-tensors-stack.tex -- the stack layout: the lattice a figure of an
% algorithm is drawn on, sites across and named layers down, and the tensors
% in its slots. It implements every operation of the abstract layout
% (tikz-tensors-layout.tex); a tree (tikz-tensors-tree.tex) is a stack that
% sizes and routes its own way.
%
% Its tensors are recorded -- the sites and layers each covers -- and that
% record is all the contraction rules below read: a tensor's indices follow
% from where it is and from its type, never from a coordinate.
\tn@kind{stack}{}

% The center of column <#2> of stack <#1>, as \tn@cx; of layer number <#2>, as
% \tn@ly. Column 0 and column n+1 are the slots just beyond the ends; a
% fraction is a line between two. Layers are a rise apart, or closer where the
% stack says so (close=), so each one's height is kept (@y@<k>); beyond the
% first and the last the stack goes on at a rise.
\def\tn@colx#1#2{\tn@len\tn@cx{\tn@get{#1@x0}+(#2)*\tn@get{#1@pitch}}}
\def\tn@layint#1#2{% the height of whole layer number <#2>, as \tn@lyi
  \pgfmathtruncatemacro\tn@lyn{#2}%
  \ifnum\tn@lyn<1
    \tn@len\tn@lyi{\tn@get{#1@y@1}+(1-\tn@lyn)*\tn@get{#1@rise}}%
  \else\ifnum\tn@lyn>\tn@get{#1@L}
    \tn@len\tn@lyi{\tn@get{#1@y@\tn@get{#1@L}}-(\tn@lyn-\tn@get{#1@L})*\tn@get{#1@rise}}%
  \else \edef\tn@lyi{\tn@get{#1@y@\tn@lyn}}\fi\fi}
\def\tn@layy#1#2{%
  \pgfmathsetmacro\tn@lyk{#2}%
  \pgfmathtruncatemacro\tn@lyf{floor(\tn@lyk)}%
  \pgfmathsetmacro\tn@lyr{\tn@lyk-\tn@lyf}%
  \tn@layint{#1}{\tn@lyf}\let\tn@ly\tn@lyi
  \ifdim\tn@lyr pt=0pt \else
    \let\tn@lyA\tn@lyi \tn@layint{#1}{\tn@lyf+1}%
    \tn@len\tn@ly{\tn@lyA+\tn@lyr*(\tn@lyi-\tn@lyA)}%
  \fi}

% ---- placing ----------------------------------------------------------------
% \tnstack[left=<label>, right=<label>, tree]{<name>}{<sites>}{<layers>}
% A grid of <sites> slots across and the <layers> down, top to bottom, placed
% after what is already on the row (an equals sign, or another stack, whose
% end slot and its own keep them a pitch apart) and centered on the row, so that
% the stacks of an equation line up. left= and right= put an environment at
% either end. Without one, the stack keeps the slot beyond its end as far as
% that slot's border, which is where an open index there ends -- so an index
% opened later never runs into the equals sign before it. [tree] makes it the
% tree layout. close= names layers that sit closer under the layer above them
% than a rise -- `closerise' (\tnset), which halves the gap between two slots
% -- as the layers of a circuit of gates are drawn. scale=, pitch=, rise=,
% size= and stub= draw it at factors of the notation's distances and sizes
% (tikz-tensors-layout.tex): a circuit of many layers, drawn compact.
% rises={<layer>=<factor>, ...} sets the gap above one layer alone, at a
% factor of the stack's rise; it is said here, where the layers' heights are
% fixed, and not in \tnlayer, because a layer's slots are named and placed
% before anything is put in them.
%
%   \tnstack[left=$L$, right=$R$]{H}{2}{ket, op, bra}
%   \tnstack[close={g2}]{S}{5}{ket, g1, g2}
%   \tnstack[scale=0.6, rise=0.8]{C}{3}{g1, g2, g3}
%   \tnstack[rises={g2=0.5, g3=0.5}]{T}{4}{ket, g1, g2, g3}
%
% The environment is the node <name>-left or <name>-right. Every slot, and the
% ones just beyond either end, is the coordinate <name>-<layer>-<i>-slot, which
% is what a label for something that is not a tensor is placed on.
\pgfkeys{/tn/stack/.is family, /tn/stack/left/.initial=,
         /tn/stack/right/.initial=, /tn/stack/close/.initial=,
         /tn/stack/rises/.initial=,
         /tn/stack/tree/.code={\def\tn@stkind{tree}}}
\pgfkeys{/tn/stack/scale/.forward to=/tn@factor/scale,
  /tn/stack/pitch/.forward to=/tn@factor/pitch,
  /tn/stack/rise/.forward to=/tn@factor/rise,
  /tn/stack/size/.forward to=/tn@factor/size,
  /tn/stack/stub/.forward to=/tn@factor/stub}
\def\tn@stack@geometry#1{\tn@set{#1@pitch}{\tn@pitch}\tn@set{#1@rise}{\tn@rise}}
% <layer>=<factor>, of rises=, as \tn@rl and \tn@rf (the factor has no `=')
\def\tn@risepair#1=#2=#3\tn@end{\edef\tn@rl{\zap@space#1 \@empty}\def\tn@rf{#2}}
\newcommand{\tnstack}[4][]{%
  \def\tn@stkind{stack}%
  \pgfkeys{/tn/stack/.cd, left=, right=, close=, rises=, #1}%
  \tn@newblock{#2}{\tn@stkind}%
  \tn@scaleblock{#2}%
  \tn@op{#2}{geometry}{#2}%
  \pgfkeysgetvalue{/tn/stack/left}\tn@envl
  \pgfkeysgetvalue{/tn/stack/right}\tn@envr
  \tn@set{#2@n}{#3}%
  \tn@set{#2@layers}{#4}%
  \global\tn@c=0
  \foreach \tn@l in {#4}{%
    \global\advance\tn@c by 1
    \tn@set{#2@k@\tn@l}{\the\tn@c}\tn@set{#2@name@\the\tn@c}{\tn@l}}%
  \tn@set{#2@L}{\the\tn@c}%
  % how far each layer is under the first, and the stack's height
  \pgfkeysgetvalue{/tn/stack/close}\tn@close
  \foreach \tn@l in \tn@close{%
    \ifcsname tn@#2@k@\tn@l\endcsname\else
      \PackageError{tikz-tensors}{close=: stack `#2' has no layer `\tn@l'}{}\fi
    \ifnum\tn@get{#2@k@\tn@l}=1
      \PackageError{tikz-tensors}{close=: `\tn@l' is the first layer of `#2';
        a layer is close to the one above it}{}\fi
    \tn@set{#2@close@\tn@l}{1}}%
  \pgfkeysgetvalue{/tn/stack/rises}\tn@rises
  \foreach \tn@rp in \tn@rises{%
    \expandafter\tn@risepair\tn@rp=\tn@end
    \ifcsname tn@#2@k@\tn@rl\endcsname\else
      \PackageError{tikz-tensors}{rises=: stack `#2' has no layer `\tn@rl'}{}\fi
    \ifnum\tn@get{#2@k@\tn@rl}=1
      \PackageError{tikz-tensors}{rises=: `\tn@rl' is the first layer of `#2';
        a rise is the gap above a layer}{}\fi
    \tn@len\tn@o{\tn@rise*(\tn@rf)}\tn@set{#2@rises@\tn@rl}{\tn@o}}%
  \tn@set{#2@off@1}{0pt}%
  \ifnum\tn@get{#2@L}>1
    \foreach \tn@k in {2,...,\tn@get{#2@L}}{%
      \edef\tn@kl{\tn@get{#2@name@\tn@k}}%
      \ifcsname tn@#2@rises@\tn@kl\endcsname \edef\tn@step{\tn@get{#2@rises@\tn@kl}}%
      \else\ifcsname tn@#2@close@\tn@kl\endcsname \let\tn@step\tn@closerise
      \else \edef\tn@step{\tn@get{#2@rise}}\fi\fi
      \tn@len\tn@o{\tn@get{#2@off@\the\numexpr\tn@k-1\relax}+\tn@step}%
      \tn@set{#2@off@\tn@k}{\tn@o}}%
  \fi
  \tn@set{#2@H}{\tn@get{#2@off@\tn@get{#2@L}}}%
  % A row after \tnbreak hangs a gap under everything drawn so far, keeping the
  % slot above its top layer as far as that slot's border -- where an index
  % opened upward ends -- as a stack keeps the slots beyond its ends.
  \tn@block{0pt}{\tn@gap+\tn@get{#2@H}/2+\tn@get{#2@rise}-\tn@slot/2}%
  \tn@set{#2@y0}{\the\tn@line@y}%
  \foreach \tn@k in {1,...,\tn@get{#2@L}}{%
    \tn@len\tn@o{\the\tn@line@y+\tn@get{#2@H}/2-\tn@get{#2@off@\tn@k}}%
    \tn@set{#2@y@\tn@k}{\tn@o}}%
  \ifx\tn@envl\@empty
    \tn@set{#2@hasl}{0}%
    \tn@len\tn@xz{\the\tn@cursor-\tn@slot/2}%
  \else
    \tn@set{#2@hasl}{1}%
    \tn@len\tn@xz{\the\tn@cursor+\tn@envw/2}%
  \fi
  \tn@set{#2@x0}{\tn@xz}%
  \tn@len\tn@h{\tn@get{#2@H}+\tn@slot}%
  \ifx\tn@envl\@empty\else
    \node[tn env, minimum height=\tn@h] (#2-left) at (\tn@xz,\the\tn@line@y) {\tn@envl};\tn@tracenode{#2-left}%
  \fi
  \ifx\tn@envr\@empty
    \tn@set{#2@hasr}{0}%
    \tn@colx{#2}{#3+1}\tn@reach{\tn@cx-\tn@slot/2}%
  \else
    \tn@set{#2@hasr}{1}%
    \tn@colx{#2}{#3+1}%
    \node[tn env, minimum height=\tn@h] (#2-right) at (\tn@cx,\the\tn@line@y) {\tn@envr};\tn@tracenode{#2-right}%
    \tn@reach{\tn@cx+\tn@envw/2}%
  \fi
  \pgfmathtruncatemacro\tn@nn{#3+1}%
  \foreach \tn@l in {#4}{%
    \tn@layy{#2}{\tn@get{#2@k@\tn@l}}%
    \foreach \tn@i in {0,...,\tn@nn}{%
      \tn@colx{#2}{\tn@i}%
      \path[overlay] (\tn@cx,\tn@ly) coordinate (#2-\tn@l-\tn@i-slot);}}%
}

% \tnlayer{<stack>}{<layer>}{<entry>, ...}
% The layer's slots, left to right, one entry each:
%   <type>/<label>                 a tensor in that slot
%   <type>/<label>/<span>          one tensor across <span> slots: a gate, a
%                                  two-site tensor, as wide as the slots
%   <type>/<label>/<span>/<down>   ... and down <down> layers, as tall as
%                                  they are, with an index into each: the
%                                  zipper of zip-up. The layers below leave
%                                  its slots as `.'
%   .                              nothing: the slot is empty
%   -                              nothing, but the layer's index runs
%                                  through, along the layer
%   |                              nothing, but the site's index runs
%                                  through, down the site
%   +                              nothing, but both run through, crossing:
%                                  a site's wire under the line of a gate
%                                  that acts on the sites either side
%   dots                           the dots of a chain that goes on (not `...'
%                                  which \foreach reads as a range)
%   <n>*<entry>                    any of these, <n> times over: a chain is
%                                  its tensor repeated, 4*canl/$A$, and a run
%                                  of empty slots is 7*.
% [size=<f>] draws the layer's tensors at <f> times the stack's size.
% A type with a comma in it goes in braces, {coef, circle}/$A$, and so does a
% label with a slash in it, gate/{$U(\delta t/2)$}/2. The tensor of
% site i is the node <stack>-<layer>-i; one that spans is named by its first
% site and its top layer.
\def\tn@entry#1/#2/#3/#4/#5\tn@end{\def\tn@sty{#1}\def\tn@lab{#2}%
  \def\tn@span{#3}\def\tn@vspan{#4}}
\def\tn@cdot{.}\def\tn@cE{-}\def\tn@cB{|}\def\tn@cX{+}\def\tn@cD{D}\def\tn@cL{L}\def\tn@cR{R}
\def\tn@cdots{dots}
% <n>*<entry> is <entry> <n> times, as \tn@n and \tn@rest. Only a whole number
% before the first `*' counts, so a label with a star in it, canl/$A^*$, is a
% label.
\def\tn@rep#1*#2\tn@end{%
  \def\tn@rest{#2}%
  \ifx\tn@rest\@empty \def\tn@n{1}\def\tn@rest{#1}%
  \else
    \edef\tn@digits{\detokenize\expandafter{\romannumeral-0#1}}%
    \ifx\tn@digits\@empty \def\tn@n{#1}\tn@repstrip#2\tn@end
    \else \def\tn@n{1}\tn@repstrip#1*#2\tn@end \fi
  \fi}
\def\tn@repstrip#1*\tn@end{\def\tn@rest{#1}}
% How wide a tensor across <span> sites is drawn, as \tn@span@w sites: as wide
% as them on a stack.
\def\tn@stack@spanwidth{\let\tn@span@w\tn@span}
\pgfkeys{/tn/layer/.is family, /tn/layer/size/.initial=1}
\newcommand{\tnlayer}[4][]{%
  \pgfkeys{/tn/layer/.cd, size=1, #1}\pgfkeysgetvalue{/tn/layer/size}\tn@lsize
  \tn@layerat{#2}{#3}{#4}}
\def\tn@layerat#1#2#3{%
  \ifcsname tn@#1@k@#2\endcsname\else
    \PackageError{tikz-tensors}{Stack `#1' has no layer `#2'}{}\fi
  \tn@tokens{#1}%
  \edef\tn@k{\tn@get{#1@k@#2}}%
  \global\tn@c=1
  \foreach \tn@each in {#3}{%
  \expandafter\tn@rep\tn@each*\tn@end
  \ifnum\tn@n<1
    \PackageError{tikz-tensors}{`\tn@n*' in layer `#2' of `#1': a run is one
      entry or more}{Leave the entry out instead of writing it 0 times.}%
    \def\tn@n{1}\def\tn@rest{.}%
  \fi
  \foreach \tn@z in {1,...,\tn@n}{%
    \expandafter\tn@entry\tn@rest/////\tn@end
    \edef\tn@i{\the\tn@c}%
    \ifx\tn@span\@empty \def\tn@span{1}\fi
    \ifx\tn@vspan\@empty \def\tn@vspan{1}\fi
    \ifx\tn@sty\tn@cdot
    \else\ifx\tn@sty\tn@cE
      \tn@set{#1@o@#2@\tn@i}{-}%
    \else\ifx\tn@sty\tn@cB
      \tn@set{#1@o@#2@\tn@i}{|}%
    \else\ifx\tn@sty\tn@cX
      \tn@set{#1@o@#2@\tn@i}{+}%
    \else\ifx\tn@sty\tn@cdots
      \tn@set{#1@o@#2@\tn@i}{D}%
      \tn@layy{#1}{\tn@k}\tn@colx{#1}{\tn@i}%
      \node[tn label] (tn@dots) at (\tn@cx,\tn@ly) {$\cdots$};\tn@tracenode{tn@dots}%
    \else
      \tn@layy{#1}{\tn@k}\let\tn@ytop\tn@ly
      \tn@layy{#1}{\tn@k+\tn@vspan-1}%
      \tn@len\tn@hh{\tn@ytop-\tn@ly+\tn@slot}%
      \tn@len\tn@ly{(\tn@ytop+\tn@ly)/2}%
      \tn@colx{#1}{\tn@i+(\tn@span-1)/2}%
      \tn@len\tn@w{(\tn@span-1)*\tn@get{#1@pitch}+\tn@slot}%
      \tn@op{#1}{spanwidth}%
      \tn@sized{#1}{\tn@lsize}%
      \edef\tn@opt{\tn@opt\ifnum\tn@span@w>1 , minimum width=\tn@w\fi
                   \ifnum\tn@vspan>1 , minimum height=\tn@hh\fi}%
      \edef\tn@nm{#1-#2-\tn@i}%
      \expandafter\tn@place\expandafter{\tn@opt}{\tn@nm}{\tn@cx}{\tn@ly}{\tn@lab}%
      \pgfmathtruncatemacro\tn@last{\tn@i+\tn@span-1}%
      \pgfmathtruncatemacro\tn@kb{\tn@k+\tn@vspan-1}%
      \tn@record{\tn@nm}{#1}{\tn@i}{\tn@last}{\tn@k}{\tn@kb}%
      \foreach \tn@q in {\tn@k,...,\tn@kb}{%
        \foreach \tn@j in {\tn@i,...,\tn@last}{%
          \tn@set{#1@o@\tn@get{#1@name@\tn@q}@\tn@j}{\tn@nm}}}%
    \fi\fi\fi\fi\fi
    \global\advance\tn@c by \tn@span\relax
  }}%
}

% What occupies slot <#3> of layer <#2> of stack <#1>, as \tn@o, and what kind
% of thing that is, as \tn@kd:
%   0 nothing  1 a tensor (\tn@o is its name)  2 `-'  3 `|'  4 `dots'
%   5 the environment on the left (column 0)  6 on the right (column n+1)
%   7 `+'
\def\tn@occ#1#2#3{%
  \ifnum#3=0 \ifnum\tn@get{#1@hasl}=1 \def\tn@o{L}\else\def\tn@o{}\fi
  \else\ifnum#3>\tn@get{#1@n} \ifnum\tn@get{#1@hasr}=1 \def\tn@o{R}\else\def\tn@o{}\fi
  \else\ifcsname tn@#1@o@#2@#3\endcsname \edef\tn@o{\tn@get{#1@o@#2@#3}}%
  \else\def\tn@o{}\fi\fi\fi
  \ifx\tn@o\@empty\def\tn@kd{0}%
  \else\ifx\tn@o\tn@cE\def\tn@kd{2}%
  \else\ifx\tn@o\tn@cB\def\tn@kd{3}%
  \else\ifx\tn@o\tn@cD\def\tn@kd{4}%
  \else\ifx\tn@o\tn@cL\def\tn@kd{5}%
  \else\ifx\tn@o\tn@cR\def\tn@kd{6}%
  \else\ifx\tn@o\tn@cX\def\tn@kd{7}%
  \else\def\tn@kd{1}\fi\fi\fi\fi\fi\fi\fi}

% From slot <#3> of site <#2>, step by <#4> through every `|' or `+' -- the site's
% index running through -- and leave the first layer that is not one in
% \tn@kw, which may be just off the grid.
\def\tn@walk#1#2#3#4{\pgfmathtruncatemacro\tn@kw{#3}\tn@walkstep{#1}{#2}{#4}}
\def\tn@walkstep#1#2#3{%
  \let\tn@next\relax
  \ifnum\tn@kw>0 \ifnum\tn@kw>\tn@get{#1@L}\else
    \edef\tn@ln{\tn@get{#1@name@\tn@kw}}%
    \tn@occ{#1}{\tn@ln}{#2}%
    \ifnum\tn@kd=7 \def\tn@kd{3}\fi
    \ifnum\tn@kd=3
      \pgfmathtruncatemacro\tn@kw{\tn@kw+(#3)}%
      \def\tn@next{\tn@walkstep{#1}{#2}{#3}}%
    \fi
  \fi\fi
  \tn@next}

% ---- named indices ------------------------------------------------------------
% The indices a type names on one side of a tensor across n sites (up, down)
% or layers (left, right) take keys in lattice units across them: the j-th of
% c at lo - 1/2 + n (j - 1/2)/c. One per site or layer, that is each site or
% layer, as a type that names none has them; one on a side the type names
% single, the middle. Any other number is spread across the tensor's own
% outline instead of the lattice (n@<t>@isp@<side>): two indices left of a
% tensor on one layer leave it a little above and below the layer's line.
\def\tn@stack@mapindex#1{%
  \tn@eachindex{#1}\tn@stackside
  \edef\tn@mone{\tn@get{n@#1@one}}%
  \foreach \tn@ms in {left,right,up,down}{%
    \tn@sideindices{#1}{\tn@ms}%
    \ifnum\tn@nonside>0
      \ifx\tn@ms\tn@sleft \def\tn@mlo{k}\def\tn@mhi{kb}\else
      \ifx\tn@ms\tn@sright \def\tn@mlo{k}\def\tn@mhi{kb}\else
        \def\tn@mlo{a}\def\tn@mhi{b}\fi\fi
      \edef\tn@mlo{\tn@get{n@#1@\tn@mlo}}\edef\tn@mhi{\tn@get{n@#1@\tn@mhi}}%
      \ifx\tn@ms\tn@mone
        \pgfmathsetmacro\tn@mlo{(\tn@mlo+\tn@mhi)/2}\let\tn@mhi\tn@mlo \fi
      \pgfmathsetmacro\tn@mn{\tn@mhi-\tn@mlo+1}%
      \pgfmathtruncatemacro\tn@msp{\tn@mn==\tn@nonside ? 0 : 1}%
      \tn@set{n@#1@isp@\tn@ms}{\tn@msp}%
      \foreach \tn@mx [count=\tn@mj] in \tn@onside{%
        \tn@ixput{#1}{\tn@mx}{\tn@ms}%
          {\tn@mlo-0.5+\tn@mn*(\tn@mj-0.5)/\tn@nonside}}%
    \fi}%
  \tn@ixorder{#1}}
\def\tn@stackside#1#2{%
  \edef\tn@sst{\noexpand\in@{,#2,}{,left,right,up,down,}}\tn@sst
  \ifin@\else
    \PackageError{tikz-tensors}{tn index: index `#1' leaves on `#2'; a tensor
      of a stack or tree has indices up, down, left or right}{}\fi}
\def\tn@isspread#1#2{\ifcsname tn@n@#1@isp@#2\endcsname
  \edef\tn@spr{\tn@get{n@#1@isp@#2}}\else \def\tn@spr{0}\fi}
% Where the port of <#1> on side <#2> with key <#3> is: its height, for one
% left or right, as \tn@ly (\tn@ixy); its column, for one up or down, as
% \tn@cx (\tn@ixx). On the lattice, or across the outline when spread.
\def\tn@ixy#1#2#3{\tn@isspread{#1}{#2}%
  \ifnum\tn@spr=1
    \tn@anc{#1}{north}\let\tn@iyt\tn@ay \tn@anc{#1}{south}%
    \tn@len\tn@ly{(\tn@iyt+\tn@ay)/2-((#3)-(\tn@get{n@#1@k}+\tn@get{n@#1@kb})/2)
      /(\tn@get{n@#1@kb}-\tn@get{n@#1@k}+1)*0.8*(\tn@iyt-\tn@ay)}%
  \else \tn@layy{\tn@get{n@#1@s}}{#3}\fi}
\def\tn@ixx#1#2#3{\tn@isspread{#1}{#2}%
  \ifnum\tn@spr=1
    \tn@anc{#1}{east}\let\tn@ixe\tn@ax \tn@anc{#1}{west}%
    \tn@len\tn@cx{(\tn@ixe+\tn@ax)/2+((#3)-(\tn@get{n@#1@a}+\tn@get{n@#1@b})/2)
      /(\tn@get{n@#1@b}-\tn@get{n@#1@a}+1)*0.8*(\tn@ixe-\tn@ax)}%
  \else \tn@colx{\tn@get{n@#1@s}}{#3}\fi}
% The keys of the ports of <#1> on side <#2> from <#3> up to (not including)
% <#4>, in order, as \tn@bk, and how many as \tn@nbk: its named indices
% there, or, for a tensor whose type names none, <#5> -- the ports the layout
% would give it there (empty if it has no index on that side).
\newcount\tn@bkc
\def\tn@bandkeys#1#2#3#4#5{%
  \gdef\tn@bk{}\global\tn@bkc=0
  \tn@isnamed{#1}%
  \ifnum\tn@named=1
    \pgfmathsetmacro\tn@bklo{#3}\pgfmathsetmacro\tn@bkhi{#4}%
    \tn@sideindices{#1}{#2}%
    \foreach \tn@bx in \tn@onside{%
      \edef\tn@bkk{\tn@get{n@#1@ik@\tn@bx}}%
      \ifdim\tn@bkk pt<\tn@bklo pt\else \ifdim\tn@bkk pt<\tn@bkhi pt
        \global\advance\tn@bkc by 1
        \ifx\tn@bk\@empty \xdef\tn@bk{\tn@bkk}\else \xdef\tn@bk{\tn@bk,\tn@bkk}\fi
      \fi\fi}%
  \else
    \tn@hasleg{#1}{#2}%
    \ifnum\tn@hasl=1 \foreach \tn@bx in {#5}{\global\advance\tn@bkc by 1
      \ifx\tn@bk\@empty \xdef\tn@bk{\tn@bx}\else \xdef\tn@bk{\tn@bk,\tn@bx}\fi}\fi
  \fi
  \edef\tn@nbk{\the\tn@bkc}}
% the <#2>th item of list <#1>, as \tn@nth
\def\tn@nthof#1#2{\foreach \tn@nx [count=\tn@nc] in #1{%
  \ifnum\tn@nc=#2 \xdef\tn@nth{\tn@nx}\fi}}

% ---- ports -------------------------------------------------------------------
% On a stack a tensor has one port up and one down per site it covers -- one,
% in the middle, on a side its type names as single -- and one left and one
% right per layer. The port up or down in column \tn@cx is in that column.
\def\tn@stack@vport#1#2{%
  \let\tn@px\tn@cx
  \tn@inner{#1}\let\tn@py\tn@ay
  \edef\tn@t{#2}%
  \ifx\tn@t\tn@sdown \tn@anc{#1}{south}\else \tn@anc{#1}{north}\fi
  \let\tn@pe\tn@ay}

% The key of the <#3>th port of tensor <#1> on side <#2>, as \tn@pkey: the
% layer it is on (left, right) or the site (up, down), or the middle of the
% sites on a side its type names single -- the key \tn@use marks it by.
\def\tn@stack@portkey#1#2#3{%
  \edef\tn@pks{#2}%
  \ifx\tn@pks\tn@sleft \pgfmathtruncatemacro\tn@pkey{\tn@get{n@#1@k}+(#3)-1}%
  \else\ifx\tn@pks\tn@sright \pgfmathtruncatemacro\tn@pkey{\tn@get{n@#1@k}+(#3)-1}%
  \else
    \edef\tn@one{\tn@get{n@#1@one}}%
    \ifx\tn@one\tn@pks \pgfmathsetmacro\tn@pkey{(\tn@get{n@#1@a}+\tn@get{n@#1@b})/2}%
    \else \pgfmathtruncatemacro\tn@pkey{\tn@get{n@#1@a}+(#3)-1}\fi
  \fi\fi}
% Every port of a tensor on a stack, side by side: one per layer it spans left
% and right, one per site up and down -- one, in the middle, on a side its
% type names single.
\def\tn@stack@ports#1{%
  \edef\tn@one{\tn@get{n@#1@one}}%
  \foreach \tn@ss in {left,right,up,down}{%
    \tn@hasleg{#1}{\tn@ss}%
    \ifnum\tn@hasl=1
      \ifx\tn@ss\tn@one
        \pgfmathsetmacro\tn@q{(\tn@get{n@#1@a}+\tn@get{n@#1@b})/2}\tn@portvisit{#1}{\tn@ss}{\tn@q}%
      \else
        \ifx\tn@ss\tn@sleft \def\tn@sr{k}\def\tn@srb{kb}\else
        \ifx\tn@ss\tn@sright \def\tn@sr{k}\def\tn@srb{kb}\else
          \def\tn@sr{a}\def\tn@srb{b}\fi\fi
        \edef\tn@slo{\tn@get{n@#1@\tn@sr}}\edef\tn@shi{\tn@get{n@#1@\tn@srb}}%
        \foreach \tn@q in {\tn@slo,...,\tn@shi}{\tn@portvisit{#1}{\tn@ss}{\tn@q}}%
      \fi
    \fi}}
\def\tn@showport#1#2#3{%
  \tn@ifused{#1}{#2}{#3}\tn@ifcut{#1}{#2}{#3}%
  \edef\tn@spd{#2}\pgfmathtruncatemacro\tn@spk{round(#3)}%
  \ifx\tn@spd\tn@sleft \def\tn@spw{on layer \tn@spk}\else
  \ifx\tn@spd\tn@sright \def\tn@spw{on layer \tn@spk}\else
  \ifx\tn@spd\tn@sup \def\tn@spw{at site \tn@spk}\else
  \ifx\tn@spd\tn@sdown \def\tn@spw{at site \tn@spk}\else \def\tn@spw{}\fi\fi\fi\fi
  \tn@fluxof{#1}{#2}{#3}%
  \typeout{\space\space index \ifx\tn@pnm\@empty #2 \tn@spw\else
    \tn@pnm\space (#2\ifx\tn@spw\@empty\else\space\tn@spw\fi)\fi
    \ifx\tn@fx\@empty\else, runs \tn@fx\ifnum\tn@fxs=0 \space(its shape)\fi\fi
    : \ifnum\tn@isused=1 drawn\else
    \ifnum\tn@iscut=1 left out (apart=)\else free\fi\fi}}

% ---- connect -----------------------------------------------------------------
% Every index the stack implies, and nothing else: along a layer between
% neighbouring slots that are both filled (or `-'), border to border; down a
% site from its first tensor to its last, through the ones between, which cover
% it; from an environment into each layer -- to the tensor there, or, if the
% slot is empty, to the border that tensor would have, which is the coordinate
% <stack>-<layer>-left or <stack>-<layer>-right (the name the end of an index
% opened there would have); and up to the dots of a chain
% that goes on. A bond with an arrow takes its direction from the shapes it
% joins. Layers named apart= have no index along them.
\def\tn@stack@connect#1{%
  \foreach \tn@l in \tn@apart{\tn@set{#1@apart@\tn@l}{1}}%
  \edef\tn@ls{\tn@get{#1@layers}}%
  \pgfmathtruncatemacro\tn@nn{\tn@get{#1@n}}%
  \foreach \tn@l in \tn@ls{%
    \edef\tn@lk{\tn@get{#1@k@\tn@l}}%
    \tn@layy{#1}{\tn@lk}\let\tn@y\tn@ly
    \ifcsname tn@#1@apart@\tn@l\endcsname\else
    \foreach \tn@i in {0,...,\tn@nn}{%
      \pgfmathtruncatemacro\tn@j{\tn@i+1}%
      \tn@occ{#1}{\tn@l}{\tn@i}\let\tn@oa\tn@o \let\tn@ka\tn@kd
      \tn@occ{#1}{\tn@l}{\tn@j}\let\tn@ob\tn@o \let\tn@kb\tn@kd
      \tn@colx{#1}{\tn@i}\let\tn@xa\tn@cx
      \tn@colx{#1}{\tn@j}\let\tn@xb\tn@cx
      \tn@hlink{\tn@along}{#1}%
    }\fi
  }%
  \ifnum\tn@nn>0
  \foreach \tn@i in {1,...,\tn@nn}{%
    \gdef\tn@prev{}%
    \foreach \tn@l in \tn@ls{%
      \tn@occ{#1}{\tn@l}{\tn@i}%
      \ifnum\tn@kd=1
        \ifx\tn@o\tn@prev\else
          \ifx\tn@prev\@empty\else \tn@vlink{#1}{\tn@prev}{\tn@o}{\tn@i}\fi
          \xdef\tn@prev{\tn@o}%
        \fi
      \fi}%
  }%
  \fi}
% The index between a tensor and the next one down its site. Two that cover
% the same sites meet once on each of them: a gate on a gate, a site on a site.
% Otherwise they meet once, in the middle of the sites they share -- a site
% under a gate, a branch under the node of a tree, an isometry of a MERA under
% the disentangler that overlaps it by half.
\def\tn@vlink#1#2#3#4{% stack, upper, lower, site
  \edef\tn@ua{\tn@get{n@#2@a}}\edef\tn@ub{\tn@get{n@#2@b}}%
  \edef\tn@na{\tn@get{n@#3@a}}\edef\tn@nb{\tn@get{n@#3@b}}%
  \tn@isnamed{#2}\let\tn@vnu\tn@named \tn@isnamed{#3}%
  \ifnum\tn@vnu\tn@named=00 \expandafter\tn@vlinkold
  \else \expandafter\tn@vpairs \fi{#1}{#2}{#3}{#4}}
% Between tensors of which one names its indices: at the first site they
% share, each index down from the upper one in the sites they share is joined
% to the index up from the lower one there that has the same place in order.
\def\tn@vpairs#1#2#3#4{%
  \pgfmathtruncatemacro\tn@voa{max(\tn@ua,\tn@na)}%
  \pgfmathtruncatemacro\tn@vob{min(\tn@ub,\tn@nb)}%
  \ifnum#4=\tn@voa
    \def\tn@vsame{0}\ifnum\tn@ua=\tn@na \ifnum\tn@ub=\tn@nb \def\tn@vsame{1}\fi\fi
    \tn@vdefault{#2}{down}\tn@bandkeys{#2}{down}{\tn@voa-0.5}{\tn@vob+0.5}{\tn@vdef}%
    \let\tn@vku\tn@bk \let\tn@vnku\tn@nbk
    \tn@vdefault{#3}{up}\tn@bandkeys{#3}{up}{\tn@voa-0.5}{\tn@vob+0.5}{\tn@vdef}%
    \let\tn@vkl\tn@bk
    \ifnum\tn@vnu=0 \ifnum\tn@vnku>\tn@nbk \let\tn@vnku\tn@nbk \fi\fi
    \tn@isnamed{#3}%
    \ifnum\tn@named=0 \ifnum\tn@nbk>\tn@vnku \let\tn@nbk\tn@vnku \fi\fi
    \ifnum\tn@vnku=\tn@nbk\else
      \PackageError{tikz-tensors}{\string\tnconnect: \tn@called{#2} has
        \tn@vnku\space \ifnum\tn@vnku=1 index\else indices\fi\space down into
        the sites it shares with \tn@called{#3}, which has \tn@nbk\space up}{The
        indices are joined in order, the first down to the first up; a type
        says them with tn index.}%
      \def\tn@nbk{0}%
    \fi
    \ifnum\tn@vnku=\tn@nbk \ifnum\tn@nbk>0
      \foreach \tn@vu [count=\tn@vj] in \tn@vku{%
        \tn@nthof\tn@vkl\tn@vj \let\tn@vl\tn@nth
        \tn@ifcut{#2}{down}{\tn@vu}\let\tn@vc\tn@iscut \tn@ifcut{#3}{up}{\tn@vl}%
        \ifnum\tn@vc\tn@iscut=00
          \tn@use{#2}{down}{\tn@vu}\tn@use{#3}{up}{\tn@vl}%
          \tn@ixx{#2}{down}{\tn@vu}\let\tn@vxu\tn@cx
          \tn@ixx{#3}{up}{\tn@vl}\let\tn@vxl\tn@cx
          \tn@orient{#2}{down}{\tn@vu}{#3}{up}{\tn@vl}{\string\tnconnect}%
          \tn@facingx{\tn@down}{#2}{#3}{\tn@vxu}{\tn@vxl}%
        \fi}%
    \fi\fi
  \fi}
% the ports a type that names no indices has toward a neighbour up or down,
% as \tn@vdef: one per shared site if the two cover the same sites, else the
% middle of the ones they share; the middle of its own on a single side
\def\tn@vdefault#1#2{%
  \edef\tn@vone{\tn@get{n@#1@one}}\def\tn@vside{#2}%
  \ifx\tn@vone\tn@vside
    \pgfmathsetmacro\tn@vdef{(\tn@get{n@#1@a}+\tn@get{n@#1@b})/2}%
  \else\ifnum\tn@vsame=1 \edef\tn@vdef{\tn@voa,...,\tn@vob}%
  \else \pgfmathsetmacro\tn@vdef{(\tn@voa+\tn@vob)/2}\fi\fi}
\def\tn@vlinkold#1#2#3#4{%
  \def\tn@draw{0}%
  \ifnum\tn@ua=\tn@na \ifnum\tn@ub=\tn@nb \def\tn@col{#4}\def\tn@draw{1}\fi\fi
  \ifnum\tn@draw=0
    \pgfmathtruncatemacro\tn@oa{max(\tn@ua,\tn@na)}%
    \pgfmathtruncatemacro\tn@ob{min(\tn@ub,\tn@nb)}%
    \ifnum#4=\tn@oa \edef\tn@col{(\tn@oa+\tn@ob)/2}\def\tn@draw{1}\fi
  \fi
  \ifnum\tn@draw=1
    \tn@hasleg{#2}{down}\ifnum\tn@hasl=0 \def\tn@draw{0}\fi
    \tn@hasleg{#3}{up}\ifnum\tn@hasl=0 \def\tn@draw{0}\fi
    \tn@ifcut{#2}{down}{\tn@col}\ifnum\tn@iscut=1 \def\tn@draw{0}\fi
    \tn@ifcut{#3}{up}{\tn@col}\ifnum\tn@iscut=1 \def\tn@draw{0}\fi
  \fi
  \ifnum\tn@draw=1
    \tn@colx{#1}{\tn@col}%
    \tn@use{#2}{down}{\tn@col}\tn@use{#3}{up}{\tn@col}%
    \tn@orient{#2}{down}{\tn@col}{#3}{up}{\tn@col}{\string\tnconnect}%
    \tn@facing{\tn@down}{#2}{#3}%
  \fi}
% From the down port of #2 to the up port of #3, under it, both in column
% \tn@cx -- or, \tn@facingx, in columns <#4> and <#5>: straight where the two
% ports line up, else down, across half way between the two borders, and down
% again.
\def\tn@facing#1#2#3{\tn@facingx{#1}{#2}{#3}{\tn@cx}{\tn@cx}}
\def\tn@facingx#1#2#3#4#5{%
  \edef\tn@fxl{#5}\edef\tn@cx{#4}%
  \tn@vport{#2}{down}\let\tn@xu\tn@px \let\tn@yu\tn@py \let\tn@eu\tn@pe
  \let\tn@cx\tn@fxl
  \tn@vport{#3}{up}\let\tn@xl\tn@px \let\tn@yl\tn@py
  \ifdim\tn@xu=\tn@xl
    \tn@dline{#1}{\tn@xu}{\tn@yu}{\tn@xl}{\tn@yl}%
  \else
    \tn@len\tn@ym{(\tn@eu+\tn@pe)/2}%
    \ifnum\tn@rev=1
      \tn@turns{#1}{(\tn@xl,\tn@yl) -- (\tn@xl,\tn@ym)
        -- (\tn@xu,\tn@ym) -- (\tn@xu,\tn@yu)}%
    \else
      \tn@turns{#1}{(\tn@xu,\tn@yu) -- (\tn@xu,\tn@ym)
        -- (\tn@xl,\tn@ym) -- (\tn@xl,\tn@yl)}%
    \fi
  \fi}
% One pair of neighbouring slots along a layer. Between two tensors of which
% one names its indices, \tn@hpairs; a tensor that names them meets anything
% else with the one index it has on the layer's line, if any.
\def\tn@hlink#1#2{%
  \def\tn@same{0}%
  \ifnum\tn@ka=1 \ifx\tn@oa\tn@ob \def\tn@same{1}\fi\fi
  \def\tn@hnm{0}%
  \ifnum\tn@same=0
    \ifnum\tn@ka=1 \tn@isnamed{\tn@oa}\ifnum\tn@named=1 \def\tn@hnm{1}\fi\fi
    \ifnum\tn@kb=1 \tn@isnamed{\tn@ob}\ifnum\tn@named=1 \def\tn@hnm{1}\fi\fi
  \fi
  \ifnum\tn@hnm=1
    \ifnum\tn@ka\tn@kb=11 \tn@hpairs{#1}%
    \else \tn@hlinkone{#1}{#2}\fi
  \else \tn@hlinkold{#1}{#2}\fi}
% Each index right of the one on this layer (its keys from lk - 1/2 to
% lk + 1/2) joined to the index left of the other that has the same place in
% order.
\def\tn@hpairs#1{%
  \tn@bandkeys{\tn@oa}{right}{\tn@lk-0.5}{\tn@lk+0.5}{\tn@lk}%
  \let\tn@hka\tn@bk \let\tn@hna\tn@nbk
  \tn@bandkeys{\tn@ob}{left}{\tn@lk-0.5}{\tn@lk+0.5}{\tn@lk}%
  % a tensor that names no indices has one wherever the other has one
  \tn@isnamed{\tn@oa}\let\tn@hma\tn@named \tn@isnamed{\tn@ob}%
  \ifnum\tn@hma=0 \ifnum\tn@hna>\tn@nbk \let\tn@hna\tn@nbk \fi\fi
  \ifnum\tn@named=0 \ifnum\tn@nbk>\tn@hna \let\tn@nbk\tn@hna \fi\fi
  \ifnum\tn@hna=\tn@nbk\else
    \PackageError{tikz-tensors}{\string\tnconnect: \tn@called{\tn@oa} has
      \tn@hna\space \ifnum\tn@hna=1 index\else indices\fi\space right on
      layer `\tn@l', and \tn@called{\tn@ob} beside it \tn@nbk\space left}{The
      indices are joined in order, the first right to the first left; a type
      says them with tn index, and apart= leaves a bond out.}%
    \def\tn@nbk{0}%
  \fi
  \ifnum\tn@hna=\tn@nbk \ifnum\tn@nbk>0
    \let\tn@hkb\tn@bk
    \foreach \tn@hu [count=\tn@hj] in \tn@hka{%
      \tn@nthof\tn@hkb\tn@hj \let\tn@hv\tn@nth
      \tn@ifcut{\tn@oa}{right}{\tn@hu}\let\tn@hc\tn@iscut
      \tn@ifcut{\tn@ob}{left}{\tn@hv}%
      \ifnum\tn@hc\tn@iscut=00
        \tn@use{\tn@oa}{right}{\tn@hu}\tn@use{\tn@ob}{left}{\tn@hv}%
        \tn@ixy{\tn@oa}{right}{\tn@hu}\let\tn@hya\tn@ly
        \tn@ixy{\tn@ob}{left}{\tn@hv}\let\tn@hyb\tn@ly
        \tn@hport{\tn@oa}{right}\let\tn@xa\tn@px
        \tn@hport{\tn@ob}{left}\let\tn@xb\tn@px
        \tn@orient{\tn@oa}{right}{\tn@hu}{\tn@ob}{left}{\tn@hv}{\string\tnconnect}%
        \tn@dline{#1}{\tn@xa}{\tn@hya}{\tn@xb}{\tn@hyb}%
      \fi}%
  \fi\fi}
% A tensor that names its indices, beside an environment, the dots or a slot
% an index runs through: it has one index on the layer's line, or none.
\def\tn@hlinkone#1#2{%
  \def\tn@hok{1}%
  \ifnum\tn@ka=1 \tn@hlinkchk{\tn@oa}{right}\tn@ka\fi
  \ifnum\tn@kb=1 \tn@hlinkchk{\tn@ob}{left}\tn@kb\fi
  \tn@hlinkold{#1}{#2}}
% none on the layer's line: as if its slot were empty
\def\tn@hlinkchk#1#2#3{%
  \tn@bandkeys{#1}{#2}{\tn@lk-0.5}{\tn@lk+0.5}{\tn@lk}%
  \ifnum\tn@nbk=0 \def#3{0}\fi
  \ifnum\tn@nbk>1
    \PackageError{tikz-tensors}{\string\tnconnect: \tn@called{#1} has
      \tn@nbk\space indices #2 on layer `\tn@l', where it meets one index}{}\fi
  \ifnum\tn@nbk=1 \ifdim\tn@bk pt=\tn@lk pt\else
    \PackageError{tikz-tensors}{\string\tnconnect: the index #2 of
      \tn@called{#1} is not on the line of layer `\tn@l', where it meets one
      index}{}\fi\fi}
\def\tn@hlinkold#1#2{%
  \ifnum\tn@same=0
  \def\tn@sa{0}\def\tn@sb{0}%
  \ifcase\tn@ka\or\def\tn@sa{1}\or\def\tn@sa{1}\or\or\or\def\tn@sa{1}\or\def\tn@sa{1}\or
    \def\tn@sa{1}\fi
  \ifcase\tn@kb\or\def\tn@sb{1}\or\def\tn@sb{1}\or\or\or\def\tn@sb{1}\or\def\tn@sb{1}\or
    \def\tn@sb{1}\fi
  % a side its type does not give it, or a bond apart= leaves out, is no bond
  \ifnum\tn@ka=1 \tn@hasleg{\tn@oa}{right}\ifnum\tn@hasl=0 \def\tn@sa{0}\fi
    \tn@ifcut{\tn@oa}{right}{\tn@lk}\ifnum\tn@iscut=1 \def\tn@sa{0}\fi\fi
  \ifnum\tn@kb=1 \tn@hasleg{\tn@ob}{left}\ifnum\tn@hasl=0 \def\tn@sb{0}\fi
    \tn@ifcut{\tn@ob}{left}{\tn@lk}\ifnum\tn@iscut=1 \def\tn@sb{0}\fi\fi
  \ifnum\tn@sa\tn@sb=11
    % a bond: from the right border of the one to the left border of the other
    \def\tn@pa{}\def\tn@pb{}%
    \ifnum\tn@ka=1
      \tn@use{\tn@oa}{right}{\tn@lk}%
      \tn@hport{\tn@oa}{right}\let\tn@xa\tn@px \let\tn@pa\tn@oa \fi
    \ifnum\tn@kb=1
      \tn@use{\tn@ob}{left}{\tn@lk}%
      \tn@hport{\tn@ob}{left}\let\tn@xb\tn@px \let\tn@pb\tn@ob \fi
    \tn@orient{\tn@pa}{right}{\tn@lk}{\tn@pb}{left}{\tn@lk}{\string\tnconnect}%
    \tn@dline{#1}{\tn@xa}{\tn@y}{\tn@xb}{\tn@y}%
  \else
    \ifnum\tn@ka=5 \ifnum\tn@sb=0 \ifnum\tn@kb=4 \else
      \tn@len\tn@xb{\tn@xb-\tn@slot/2}%
      \tn@line{#1}{\tn@xa}{\tn@y}{\tn@xb}{\tn@y}%
      \coordinate (#2-\tn@l-left) at (\tn@xb,\tn@y);
      \tn@reach{\tn@xb}%
    \fi\fi\fi
    \ifnum\tn@kb=6 \ifnum\tn@sa=0 \ifnum\tn@ka=4 \else
      \tn@len\tn@xa{\tn@xa+\tn@slot/2}%
      \tn@line{#1}{\tn@xb}{\tn@y}{\tn@xa}{\tn@y}%
      \coordinate (#2-\tn@l-right) at (\tn@xa,\tn@y);
    \fi\fi\fi
    \ifnum\tn@ka\tn@kb=14
      \tn@use{\tn@oa}{right}{\tn@lk}%
      \tn@hport{\tn@oa}{right}\let\tn@xa\tn@px \tn@len\tn@xb{\tn@xb-\tn@slot/2}%
      \tn@orient{\tn@oa}{right}{\tn@lk}{}{}{}{\string\tnconnect}%
      \tn@dline{#1}{\tn@xa}{\tn@y}{\tn@xb}{\tn@y}%
    \fi
    \ifnum\tn@ka\tn@kb=41
      \tn@use{\tn@ob}{left}{\tn@lk}%
      \tn@hport{\tn@ob}{left}\let\tn@xb\tn@px \tn@len\tn@xa{\tn@xa+\tn@slot/2}%
      \tn@orient{}{}{}{\tn@ob}{left}{\tn@lk}{\string\tnconnect}%
      \tn@dline{#1}{\tn@xa}{\tn@y}{\tn@xb}{\tn@y}%
    \fi
  \fi\fi}

% ---- open ------------------------------------------------------------------------
% An open index ends at the border the tensor in the next slot would have --
% past any `|', where the site's index runs on -- and shows at least `stub'
% beyond its own tensor. Opened on a tensor: its own indices that way, one per
% site it spans up or down and one per layer it spans left or right, ending at
% <tensor>-<direction>, or -<direction>-1, -2, ... when there are several. On
% the stack: every site's index from its lowest tensor (down) or highest (up),
% ending at <stack>-<i>-down / -up; or every layer's from its last tensor
% (right) or first (left), ending at <stack>-<layer>-right / -left. An
% open index runs the way its tensor says (tn flux, tn index): into it, drawn
% from the end to the tensor, or out of it, or, said neither way, outward.
%
% one index across, from tensor #2 on layer number #3, ending at coordinate #4
\def\tn@hopen#1#2#3#4{%
  \edef\tn@s{\tn@get{n@#2@s}}\tn@fluxof{#2}{\tn@dir}{#3}%
  \ifx\tn@fx\tn@sin \def\tn@rev{1}\else \def\tn@rev{0}\fi
  \tn@use{#2}{\tn@dir}{#3}%
  \tn@ixy{#2}{\tn@dir}{#3}\let\tn@hy\tn@ly
  \tn@hport{#2}{\tn@dir}\let\tn@xa\tn@px
  \ifx\tn@dir\tn@sright
    \tn@len\tn@xc{\tn@pe+\tn@stub}%
    \tn@colx{\tn@s}{\tn@get{n@#2@b}+1}\tn@len\tn@xb{\tn@cx-\tn@slot/2}%
    \ifdim\tn@xc>\tn@xb \let\tn@xb\tn@xc \fi
    \tn@reach{\tn@xb}%
    \tn@dline{#1}{\tn@xa}{\tn@hy}{\tn@xb}{\tn@hy}%
  \else
    \tn@len\tn@xc{\tn@pe-\tn@stub}%
    \tn@colx{\tn@s}{\tn@get{n@#2@a}-1}\tn@len\tn@xb{\tn@cx+\tn@slot/2}%
    \ifdim\tn@xc<\tn@xb \let\tn@xb\tn@xc \fi
    \tn@dline{#1}{\tn@xa}{\tn@hy}{\tn@xb}{\tn@hy}%
  \fi
  \coordinate (#4) at (\tn@xb,\tn@hy);
  \tn@openlabel{#4}}
% one index up or down, from tensor #2 on site #3, ending at coordinate #4;
% \tn@vopenx draws it at the column <x> instead, walking site #3. It ends at the
% border of the next slot that is not `|'; where the tensor's port is not over
% that slot (in a tree), it turns a corner half way there.
\def\tn@vopen#1#2#3#4{\tn@vopenx{#1}{#2}{#3}{#3}{#4}}
\def\tn@vopenx#1#2#3#4#5{%
  \edef\tn@s{\tn@get{n@#2@s}}%
  \tn@use{#2}{\tn@dir}{#4}%
  \ifx\tn@dir\tn@sdown
    \tn@walk{\tn@s}{#3}{\tn@get{n@#2@kb}+1}{1}%
    \tn@layy{\tn@s}{\tn@kw}\tn@len\tn@yb{\tn@ly+\tn@slot/2}%
  \else
    \tn@walk{\tn@s}{#3}{\tn@get{n@#2@k}-1}{-1}%
    \tn@layy{\tn@s}{\tn@kw}\tn@len\tn@yb{\tn@ly-\tn@slot/2}%
  \fi
  \tn@ixx{#2}{\tn@dir}{#4}%
  \tn@vport{#2}{\tn@dir}%
  \tn@fluxof{#2}{\tn@dir}{#4}%
  \ifx\tn@fx\tn@sin \def\tn@rev{1}\else \def\tn@rev{0}\fi
  \ifdim\tn@px=\tn@cx
    \tn@dline{#1}{\tn@cx}{\tn@py}{\tn@cx}{\tn@yb}%
  \else
    \tn@len\tn@ym{(\tn@pe+\tn@yb)/2}%
    \ifnum\tn@rev=1
      \tn@turns{#1}{(\tn@cx,\tn@yb) -- (\tn@cx,\tn@ym)
        -- (\tn@px,\tn@ym) -- (\tn@px,\tn@py)}%
    \else
      \tn@turns{#1}{(\tn@px,\tn@py) -- (\tn@px,\tn@ym)
        -- (\tn@cx,\tn@ym) -- (\tn@cx,\tn@yb)}%
    \fi
  \fi
  \coordinate (#5) at (\tn@cx,\tn@yb);
  \tn@openlabel{#5}}
% One port of tensor #2, side #3, key #4, opened that way: ending at
% <tensor>-<side> if it is the only one on that side, else -<side>-<m>, the
% <m>th of them -- the names \tnopen gives on a tensor.
\def\tn@stack@openport#1#2#3#4{%
  \edef\tn@dir{#3}\tn@vdir
  \ifnum\tn@isv=1 \edef\tn@lo{\tn@get{n@#2@a}}\edef\tn@hi{\tn@get{n@#2@b}}%
  \else \edef\tn@lo{\tn@get{n@#2@k}}\edef\tn@hi{\tn@get{n@#2@kb}}\fi
  \edef\tn@one{\tn@get{n@#2@one}}%
  \ifx\tn@one\tn@dir \let\tn@hi\tn@lo \fi
  \ifx\tn@pnm\@empty
    \ifnum\tn@lo=\tn@hi \def\tn@onm{#2-#3}%
    \else \pgfmathtruncatemacro\tn@m{(#4)-\tn@lo+1}\edef\tn@onm{#2-#3-\tn@m}\fi
  \else \edef\tn@onm{#2-\tn@pnm}\fi
  \ifnum\tn@isv=1
    \ifx\tn@one\tn@dir \tn@vopenx{#1}{#2}{\tn@lo}{#4}{\tn@onm}%
    \else \pgfmathtruncatemacro\tn@osite{round(#4)}%
      \tn@vopenx{#1}{#2}{\tn@osite}{#4}{\tn@onm}\fi
  \else \tn@hopen{#1}{#2}{#4}{\tn@onm}\fi}
\def\tn@vdir{\def\tn@isv{1}\ifx\tn@dir\tn@sleft\def\tn@isv{0}\fi
  \ifx\tn@dir\tn@sright\def\tn@isv{0}\fi}
\def\tn@stack@opennode#1#2{%
  \tn@needleg{#2}{\tn@dir}{\string\tnopen}%
  \tn@isnamed{#2}%
  \ifnum\tn@named=1 \expandafter\tn@opennamed \else \expandafter\tn@opennodeold \fi
  {#1}{#2}}
% A tensor that names its indices: each of them on that side, ending at
% <tensor>-<name>.
\def\tn@opennamed#1#2{%
  \edef\tn@ond{\tn@dir}\tn@sideindices{#2}{\tn@dir}%
  \foreach \tn@pnm in \tn@onside{%
    \edef\tn@onk{\tn@get{n@#2@ik@\tn@pnm}}%
    \tn@free{#2}{\tn@ond}{\tn@onk}{\string\tnopen}%
    \tn@stack@openport{#1}{#2}{\tn@ond}{\tn@onk}}}
\def\tn@opennodeold#1#2{%
  \tn@vdir
  \ifnum\tn@isv=1
    \edef\tn@a{\tn@get{n@#2@a}}\edef\tn@b{\tn@get{n@#2@b}}%
    \edef\tn@one{\tn@get{n@#2@one}}%
    \ifx\tn@one\tn@dir
      \tn@free{#2}{\tn@dir}{(\tn@a+\tn@b)/2}{\string\tnopen}%
      \tn@vopenx{#1}{#2}{\tn@a}{(\tn@a+\tn@b)/2}{#2-\tn@dir}%
    \else
    \foreach \tn@i [count=\tn@m] in {\tn@a,...,\tn@b}{%
      \tn@free{#2}{\tn@dir}{\tn@i}{\string\tnopen}%
      \ifnum\tn@a=\tn@b \tn@vopen{#1}{#2}{\tn@i}{#2-\tn@dir}%
      \else \tn@vopen{#1}{#2}{\tn@i}{#2-\tn@dir-\tn@m}\fi}%
    \fi
  \else
    \edef\tn@ka{\tn@get{n@#2@k}}\edef\tn@kb{\tn@get{n@#2@kb}}%
    \foreach \tn@q [count=\tn@m] in {\tn@ka,...,\tn@kb}{%
      \tn@free{#2}{\tn@dir}{\tn@q}{\string\tnopen}%
      \ifnum\tn@ka=\tn@kb \tn@hopen{#1}{#2}{\tn@q}{#2-\tn@dir}%
      \else \tn@hopen{#1}{#2}{\tn@q}{#2-\tn@dir-\tn@m}\fi}%
  \fi}
\def\tn@stack@openblock#1#2{%
  \tn@vdir
  \edef\tn@ls{\tn@get{#2@layers}}%
  \pgfmathtruncatemacro\tn@nn{\tn@get{#2@n}}%
  \ifnum\tn@isv=1
    \foreach \tn@i in {1,...,\tn@nn}{%
      \gdef\tn@first{}\gdef\tn@lastn{}%
      \foreach \tn@l in \tn@ls{%
        \tn@occ{#2}{\tn@l}{\tn@i}%
        \ifnum\tn@kd=1
          \ifx\tn@first\@empty \xdef\tn@first{\tn@o}\fi
          \xdef\tn@lastn{\tn@o}%
        \fi}%
      \ifx\tn@dir\tn@sdown \let\tn@pick\tn@lastn \else \let\tn@pick\tn@first \fi
      \ifx\tn@pick\@empty\else \tn@hasleg{\tn@pick}{\tn@dir}\ifnum\tn@hasl=0
        \let\tn@pick\@empty\fi\fi
      \ifx\tn@pick\@empty\else
        \edef\tn@one{\tn@get{n@\tn@pick @one}}%
        \tn@isnamed{\tn@pick}%
        \ifnum\tn@named=1
          \tn@openband{#1}{\tn@pick}{\tn@i}%
        \else\ifx\tn@one\tn@dir
          \edef\tn@a{\tn@get{n@\tn@pick @a}}\edef\tn@b{\tn@get{n@\tn@pick @b}}%
          \ifnum\tn@i=\tn@a
            \tn@ifused{\tn@pick}{\tn@dir}{(\tn@a+\tn@b)/2}%
            \ifnum\tn@isused=0
              \tn@vopenx{#1}{\tn@pick}{\tn@i}{(\tn@a+\tn@b)/2}{#2-\tn@i-\tn@dir}%
            \fi
          \fi
        \else
          \tn@ifused{\tn@pick}{\tn@dir}{\tn@i}%
          \ifnum\tn@isused=0 \tn@vopen{#1}{\tn@pick}{\tn@i}{#2-\tn@i-\tn@dir}\fi
        \fi\fi\fi}%
  \else
    \foreach \tn@l in \tn@ls{%
      \gdef\tn@first{}\gdef\tn@lastn{}%
      \foreach \tn@i in {1,...,\tn@nn}{%
        \tn@occ{#2}{\tn@l}{\tn@i}%
        \ifnum\tn@kd=1
          \ifx\tn@first\@empty \xdef\tn@first{\tn@o}\fi
          \xdef\tn@lastn{\tn@o}%
        \fi}%
      \ifx\tn@dir\tn@sright \let\tn@pick\tn@lastn \else \let\tn@pick\tn@first \fi
      \ifx\tn@pick\@empty\else \tn@hasleg{\tn@pick}{\tn@dir}\ifnum\tn@hasl=0
        \let\tn@pick\@empty\fi\fi
      \ifx\tn@pick\@empty\else
        \edef\tn@lk{\tn@get{#2@k@\tn@l}}%
        \tn@isnamed{\tn@pick}%
        \ifnum\tn@named=1 \tn@openband{#1}{\tn@pick}{\tn@lk}%
        \else
          \tn@ifused{\tn@pick}{\tn@dir}{\tn@lk}%
          \ifnum\tn@isused=0 \tn@hopen{#1}{\tn@pick}{\tn@lk}{#2-\tn@l-\tn@dir}\fi
        \fi
      \fi}%
  \fi}

% the free named indices of <#2> toward \tn@dir at site or layer <#3>
\def\tn@openband#1#2#3{%
  \edef\tn@obd{\tn@dir}\tn@sideindices{#2}{\tn@dir}%
  \foreach \tn@pnm in \tn@onside{%
    \edef\tn@onk{\tn@get{n@#2@ik@\tn@pnm}}%
    \ifdim\tn@onk pt<\dimexpr #3pt-0.5pt\relax\else
    \ifdim\tn@onk pt<\dimexpr #3pt+0.5pt\relax
      \tn@ifused{#2}{\tn@obd}{\tn@onk}%
      \ifnum\tn@isused=0 \tn@stack@openport{#1}{#2}{\tn@obd}{\tn@onk}\fi
    \fi\fi}}

% The two indices down from a tensor across two sites, exchanged on the way,
% drawn by \tn@swap and ending where \tnopen{down} would end them
% (<tensor>-down-1 and -2, by where they end). The crossing needs room, which
% a layer of `|' under the tensor gives it.
\def\tn@stack@openswap#1#2{%
  \edef\tn@s{\tn@get{n@#2@s}}\edef\tn@a{\tn@get{n@#2@a}}%
  \tn@free{#2}{down}{\tn@a}{\string\tnopenswap}\tn@use{#2}{down}{\tn@a}%
  \tn@free{#2}{down}{\tn@a+1}{\string\tnopenswap}\tn@use{#2}{down}{\tn@a+1}%
  \tn@walk{\tn@s}{\tn@a}{\tn@get{n@#2@kb}+1}{1}%
  \tn@layy{\tn@s}{\tn@kw}\tn@len\tn@yb{\tn@ly+\tn@slot/2}%
  \tn@anc{#2}{south}\tn@d=\tn@ay\relax\edef\tn@ys{\the\tn@d}%
  \tn@len\tn@drop{\tn@ys-\tn@yb}%
  \tn@colx{\tn@s}{\tn@a}\let\tn@xa\tn@cx
  \tn@colx{\tn@s}{\tn@a+1}\let\tn@xb\tn@cx
  \tn@swap{#1}{(\tn@xa,\tn@ys)}{(\tn@xb,\tn@ys)}{\tn@drop}%
  \coordinate (#2-down-1) at (\tn@xa,\tn@yb);
  \coordinate (#2-down-2) at (\tn@xb,\tn@yb);
  \tn@openlabel{#2-down-1}\tn@openlabel{#2-down-2}}

% ---- join --------------------------------------------------------------------
% <side> is up, down, left or right; <n> is the <n>th site the tensor spans
% (up, down) or layer (left, right), and the middle on a side its type names
% as single. There is no route to choose, and none goes through a tensor. Two
% ports that face each other on one line, with nothing in the slots between
% them, are joined straight. Any other index leaves each port by half a step,
% into the gutter beside it -- the line half way between two columns, or two
% layers, where no tensor sits -- and runs along gutters from one to the
% other: along the one they share, else by one gutter across (between their
% layers, or their columns), else, for two ports on one layer that do not
% face each other or have tensors between them, round the top of the stack.
% Only a tensor drawn across a gutter (one spanning several sites or layers)
% can be in the way.
%
% Port <#1> as \tn@j@<#2>x, y (where the index starts), gx, gy (where it
% reaches the gutter), d (its side), h (1 if it leaves sideways), c (its
% column, for an index up or down, or layer, for one sideways, in lattice
% units), lo and hi (the sites it spans, or the layers) and s (its stack).
\def\tn@jset#1#2{\expandafter\xdef\csname tn@j@#1\endcsname{#2}}
\def\tn@jport#1#2#3#4{%
  \edef\tn@pn{#1}\edef\tn@ps{#2}\edef\tn@q{#3}\def\tn@jp{#4}%
  \edef\tn@s{\tn@get{n@\tn@pn @s}}\tn@jset{#4s}{\tn@s}%
  \tn@jset{#4d}{\tn@ps}%
  \ifx\tn@ps\tn@sleft \def\tn@ph{1}\else\ifx\tn@ps\tn@sright \def\tn@ph{1}%
  \else \def\tn@ph{0}\fi\fi
  \tn@jset{#4h}{\tn@ph}%
  \edef\tn@a{\tn@get{n@\tn@pn @a}}\edef\tn@b{\tn@get{n@\tn@pn @b}}%
  \edef\tn@ka{\tn@get{n@\tn@pn @k}}\edef\tn@kb{\tn@get{n@\tn@pn @kb}}%
  \ifnum\tn@ph=1
    \tn@jset{#4lo}{\tn@a}\tn@jset{#4hi}{\tn@b}%
    \pgfmathtruncatemacro\tn@qc{round(\tn@q)}\tn@jset{#4c}{\tn@qc}%
    \tn@free{\tn@pn}{\tn@ps}{\tn@q}{\string\tnjoin}\tn@use{\tn@pn}{\tn@ps}{\tn@q}%
    \tn@ixy{\tn@pn}{\tn@ps}{\tn@q}%
    \tn@hport{\tn@pn}{\tn@ps}%
    \ifx\tn@ps\tn@sright \tn@colx{\tn@s}{\tn@b+0.5}\else \tn@colx{\tn@s}{\tn@a-0.5}\fi
    \tn@jset{#4x}{\tn@px}\tn@jset{#4y}{\tn@ly}%
    \tn@jset{#4gx}{\tn@cx}\tn@jset{#4gy}{\tn@ly}%
  \else
    \tn@jset{#4lo}{\tn@ka}\tn@jset{#4hi}{\tn@kb}%
    \tn@jset{#4c}{\tn@q}%
    \tn@free{\tn@pn}{\tn@ps}{\tn@q}{\string\tnjoin}\tn@use{\tn@pn}{\tn@ps}{\tn@q}%
    \tn@ixx{\tn@pn}{\tn@ps}{\tn@q}%
    \tn@vport{\tn@pn}{\tn@ps}%
    \ifx\tn@ps\tn@sdown \tn@layy{\tn@s}{\tn@kb+0.5}\else \tn@layy{\tn@s}{\tn@ka-0.5}\fi
    \tn@jset{#4x}{\tn@px}\tn@jset{#4y}{\tn@py}%
    \tn@jset{#4gx}{\tn@px}\tn@jset{#4gy}{\tn@ly}%
  \fi}
% \tn@jclear{<stack>}{<layer|site>}{<from>}{<to>}{<h>}: is every slot strictly
% between <from> and <to> empty -- along layer number <layer> (h = 1) or down
% site <site> (h = 0)? As \tn@clear.
\def\tn@jclear#1#2#3#4#5{%
  \def\tn@clear{1}%
  \pgfmathtruncatemacro\tn@lo{min(#3,#4)+1}\pgfmathtruncatemacro\tn@hi{max(#3,#4)-1}%
  \ifnum\tn@lo>\tn@hi\else
    \foreach \tn@w in {\tn@lo,...,\tn@hi}{%
      \ifnum#5=1 \tn@occ{#1}{\tn@get{#1@name@#2}}{\tn@w}%
      \else \tn@occ{#1}{\tn@get{#1@name@\tn@w}}{#2}\fi
      \ifnum\tn@kd=0 \else \gdef\tn@clear{0}\fi}%
  \fi}
\def\tn@stack@join#1#2#3#4#5#6#7{%
  \tn@jport{#2}{#3}{#4}{a}\tn@jport{#5}{#6}{#7}{b}%
  \tn@orient{#2}{#3}{#4}{#5}{#6}{#7}{\string\tnjoin}%
  \edef\tn@s{\tn@j@as}%
  % straight: facing on one line, nothing between
  \def\tn@face{0}%
  \ifnum\tn@j@ah\tn@j@bh=11 \ifdim\tn@j@ay=\tn@j@by
    \ifx\tn@j@ad\tn@sright \ifx\tn@j@bd\tn@sleft
      \ifdim\tn@j@ax<\tn@j@bx \def\tn@face{1}\fi\fi\fi
    \ifx\tn@j@ad\tn@sleft \ifx\tn@j@bd\tn@sright
      \ifdim\tn@j@ax>\tn@j@bx \def\tn@face{1}\fi\fi\fi
    \ifnum\tn@face=1
      \ifx\tn@j@ad\tn@sright \tn@jclear{\tn@s}{\tn@j@ac}{\tn@j@ahi}{\tn@j@blo}{1}%
      \else \tn@jclear{\tn@s}{\tn@j@ac}{\tn@j@bhi}{\tn@j@alo}{1}\fi
      \ifnum\tn@clear=0 \def\tn@face{0}\fi
    \fi
  \fi\fi
  \ifnum\tn@j@ah\tn@j@bh=0 \ifdim\tn@j@ax=\tn@j@bx
    \ifx\tn@j@ad\tn@sdown \ifx\tn@j@bd\tn@sup
      \ifdim\tn@j@ay>\tn@j@by \def\tn@face{1}\fi\fi\fi
    \ifx\tn@j@ad\tn@sup \ifx\tn@j@bd\tn@sdown
      \ifdim\tn@j@ay<\tn@j@by \def\tn@face{1}\fi\fi\fi
    \ifnum\tn@face=1
      \pgfmathtruncatemacro\tn@site{\tn@j@ac}%
      \ifx\tn@j@ad\tn@sdown \tn@jclear{\tn@s}{\tn@site}{\tn@j@ahi}{\tn@j@blo}{0}%
      \else \tn@jclear{\tn@s}{\tn@site}{\tn@j@bhi}{\tn@j@alo}{0}\fi
      \ifnum\tn@clear=0 \def\tn@face{0}\fi
    \fi
  \fi\fi
  \ifnum\tn@face=1
    \tn@dline{#1}{\tn@j@ax}{\tn@j@ay}{\tn@j@bx}{\tn@j@by}%
  \else
    % along the gutters: from (gx, gy) of the one to (gx, gy) of the other
    \ifnum\tn@j@ah\tn@j@bh=11
      \ifdim\tn@j@agx=\tn@j@bgx \def\tn@via{}\def\tn@viar{}%
      \else
        \ifdim\tn@j@ay=\tn@j@by
          \tn@layy{\tn@s}{0.5}% round the top of the stack
        \else
          \tn@layy{\tn@s}{min(\tn@j@ac,\tn@j@bc)+0.5}% between their layers
        \fi
        \edef\tn@via{-- (\tn@j@agx,\tn@ly) -- (\tn@j@bgx,\tn@ly)}%
        \edef\tn@viar{-- (\tn@j@bgx,\tn@ly) -- (\tn@j@agx,\tn@ly)}%
      \fi
    \else\ifnum\tn@j@ah\tn@j@bh=0
      \ifdim\tn@j@agy=\tn@j@bgy \def\tn@via{}\def\tn@viar{}%
      \else
        \ifdim\tn@j@ax=\tn@j@bx
          \tn@colx{\tn@s}{\tn@j@ac+0.5}% beside them
        \else
          \tn@colx{\tn@s}{min(\tn@j@ac,\tn@j@bc)+0.5}% between their columns
        \fi
        \tn@reach{\tn@cx}%
        \edef\tn@via{-- (\tn@cx,\tn@j@agy) -- (\tn@cx,\tn@j@bgy)}%
        \edef\tn@viar{-- (\tn@cx,\tn@j@bgy) -- (\tn@cx,\tn@j@agy)}%
      \fi
    \else\ifnum\tn@j@ah=1
      \edef\tn@via{-- (\tn@j@agx,\tn@j@bgy)}\let\tn@viar\tn@via
    \else
      \edef\tn@via{-- (\tn@j@bgx,\tn@j@agy)}\let\tn@viar\tn@via
    \fi\fi\fi
    \ifnum\tn@rev=1
      \tn@turns{#1}{(\tn@j@bx,\tn@j@by) -- (\tn@j@bgx,\tn@j@bgy) \tn@viar
        -- (\tn@j@agx,\tn@j@agy) -- (\tn@j@ax,\tn@j@ay)}%
    \else
      \tn@turns{#1}{(\tn@j@ax,\tn@j@ay) -- (\tn@j@agx,\tn@j@agy) \tn@via
        -- (\tn@j@bgx,\tn@j@bgy) -- (\tn@j@bx,\tn@j@by)}%
    \fi
  \fi}
